\documentclass{beamer}

\mode<presentation>
\usetheme{Warsaw}
\usefonttheme[onlylarge]{structurebold}
\usecolortheme[RGB={255,0,0}]{structure}
\setbeamerfont{frametitle}{size=\normalsize,series=\bfseries}
\setbeamertemplate{navigation symbols}{}
%\setbeamertemplate{background canvas}{bg=gray}
\setbeamertemplate{background canvas}{\includegraphics[width=\paperwidth,height=\paperheight]{images/openlogo-nd.eps}}
\setbeamercovered{transparent=95}
\setbeamertemplate{blocks}[rounded][shadow=true]
%\setbeamersize{margin left=0.5cm,margin right=-1cm}



\usepackage[brazil]{babel}
\usepackage[utf8]{inputenc}
\usepackage{times}
\usepackage[T1]{fontenc}

\usepackage{tikz}
\usetikzlibrary{arrows}
\tikzstyle{block}=[draw opacity=0.7,line width=1.4cm]

\title[Sysadmin: O que você é e como usar o Debian para ser um sysadmin melhor]
{%
   Sysadmin: O que você é e como usar o Debian para ser um sysadmin melhor%
}

\author[Fernando Ike de Oliveira]
{
   \textcolor{black}{Fernando Ike de Oliveira}%\inst{1}
}

\institute[Debian-BR]
%{
% \inst{1}%
% \textcolor{black}{Debian}
%}

\date[Short Occasion]
{1º /dev/net}
%\logo{\includegraphics[scale=0.1]{images/b2br}}
%\logo{\includegraphics[scale=0.1,center]{images/pgbr.eps}}

\begin{document}
\begin{frame}
\titlepage
\end{frame}

\begin{frame}
    \frametitle{O que é um sysadmin?}
    \begin{block}{É um...}
        \begin{itemize}
         \item Analista de suporte?
         \item Analista de sistema?
         \item Analista de rede?
         \item Desenvolvedor?
         \item DBA?
         \item Operador de backup?
         \item Técnico de Hardware?
	 \end{itemize}
    \end{block}
\end{frame}


\begin{frame}
    \frametitle{Definição}
        \begin{block}{Um sysadmin segundo wikipedia é...}
	 \begin{itemize}
	    \item É um profissional que mantém e opera um systema de computação e rede. Usualmente ele instala, mantém e administra servidores. Dentre suas tarefas inclui programação (scripting), gerenciamento de projeto e de operadores.
	    \item Andrelop diz: Só idiotas podem fazer o que fazemos e ainda gosta dele, por isso, nós somos malucos
         \end{itemize}
    \end{block}
\end{frame}

\begin{frame}
    \frametitle{O que realmente o sysadmin faz...}
        \begin{block}{ }
	 \begin{itemize}
	    \item É o Suporte Técnico
	    \item É responsável pelo backup
	    \item É responsável pela rede
	    \item É responsável pelo banco de dados
	    \item É responsável segurança
         \end{itemize}
    \end{block}
\end{frame}

\begin{frame}{Coisas que não são boas\dots}
  \begin{center}
     \pgfputat{\pgfxy(-113,-100)}{\pgfimage[width=\paperwidth,height=\paperheight]{images/donttouch.eps}}
         \par
  \end{center}
\end{frame}


\begin{frame}{Coisas que não são boas\dots}
  \begin{center}
     \pgfputat{\pgfxy(-113,-100)}{\pgfimage[width=\paperwidth,height=\paperheight]{images/before.eps}}
         \par
  \end{center}
\end{frame}


\begin{frame}
    \frametitle{Coisas de um péssimo sysadmin}
        \begin{block}{Nunca faça ou seja}
	 \begin{itemize}
	    \item Um fanboy
	    \item Reativo
	    \item Não ser organizado
	    \item Não documentar
         \end{itemize}
    \end{block}
\end{frame}

\begin{frame}
    \frametitle{Um bom sysadmin}
        \begin{block}{Boas práticas}
	 \begin{itemize}
	    \item Documentar
	    \item Cria processos
	    \item Não ter sempre a razão
	    \item Seu usuário é seu principal cliente
         \end{itemize}
    \end{block}
\end{frame}



\begin{frame}{Se for um bom sysadmin\dots}
  \begin{center}
     \pgfputat{\pgfxy(-113,-100)}{\pgfimage[width=\paperwidth,height=\paperheight]{images/after.eps}}
         \par
  \end{center}
\end{frame}



\begin{frame}
    \frametitle{Argumentos para convencer seu chefe usar Debian}
        \begin{block}{Convença seu chefe}
	 \begin{itemize}
	    \item ciclo de lançamento
	    \item estabalidade
	    \item atualizações de segurança não quebram os pacotes
	    \item custo inicial baixo
	    \item suporte
         \end{itemize}
    \end{block}
\end{frame}

\begin{frame}
    \frametitle{Argumentos para convencer seus colegas de trabalho}
        \begin{block}{Convença seus amigos}
	 \begin{itemize}
	    \item Pela preguiça (repetição de trabalhos chatos)
	    \item Pela cobiça (ajudar ao projeto)
         \end{itemize}
    \end{block}
\end{frame}


\begin{frame}
    \frametitle{Projetos criados pelo Debian...}
        \begin{block}{... e copiados pelos projetos }
	 \begin{itemize}
	    \item alternatives
	    \item menu
	    \item comunidade
	    \item pacotes e dependência
	    \item DFSG
         \end{itemize}
    \end{block}
\end{frame}

\begin{frame}
    \frametitle{coisas que o Debian tem para você}
        \begin{block}{ }
	 \begin{itemize}
	    \item cron-apt 
%	    \item unattended-upgrades
%	    \item logbook com dch 
%	    \item logwatch
	    \item apt pinning 
	    \item diversion
	    \item alternative
	    \item debconf 
	    \item debsums
	    \item apache2
	    \item postgresql
            \item FAI
	    \item Debian-Installer
        \end{itemize}
    \end{block}
\end{frame}

\begin{frame}
    \frametitle{coisas que o Debian tem para você}
        \begin{block}{ }
	 \begin{itemize}
 	   \item CFengine
	   \item debmirror
	   \item live-helper
	   \item equivs
	   \item statoverride
	   \item selinux
	   \item pax/grsecurity
	   \item Virtualização (XEN, Qemu, KVM, Vserver, Virtualbox, Bochs)
        \end{itemize}
    \end{block}
\end{frame}



\begin{frame}
    \frametitle{onde vc tem que ler}
        \begin{block}{leitura obrigatória}
	 \begin{itemize}
	    \item lista debian-announce
	    \item lista debian-news
	    \item lista debian-security
	    \item site bugs.debian.org
	    \item changelog.Debian.gz, README.Debian.gz
         \end{itemize}
    \end{block}
\end{frame}

\begin{frame}
    \frametitle{Mantra dos pacotes}
        \begin{block}{boas práticas}
	 \begin{itemize}
	    \item nunca misture os repositorios
	    \item nao use apt-get
	    \item aptitude é seu amigo 
         \end{itemize}
    \end{block}
\end{frame}

\begin{frame}
    \frametitle{o que é bom saber}
        \begin{block}{não deixe de aprender}
	 \begin{itemize}
            \item linguagem de programação (Perl, Python, Ruby)
	    \item ferramentas de depuração e trace
            \item resolução de gargalos e performence
         \end{itemize}
    \end{block}
\end{frame}


\begin{frame}
    \frametitle{\textbf{Referência}}
    \begin{block}{ }
        \begin{itemize}
            \item {\textit{\url{http://lists.debian.org}}}
            \item {\textit{\url{http://wiki.debian.org}}}
            \item {\textit{\url{http://www.debian.org/doc}}}
        \end{itemize}
    \end{block}

\end{frame}





\begin{frame}
    \frametitle{\textbf{Conclusão}}
    \begin{block}{Contatos:}
        \begin{itemize}
            \item {\texttt{fernando.ike@b2br.com.br}}
            \item {\texttt{fernando.ike@gmail.com}}
            \item {\textit{\url{ http://www.midstorm.org/~fike/weblog}}}
        \end{itemize}
    \end{block}

\end{frame}

\begin{frame}
    \frametitle{\textbf{Conclusão}}
    \begin{exampleblock}{Contatos:}
        \begin{itemize}
            \item {\texttt{fernando.ike@b2br.com.br}}
            \item {\texttt{fernando.ike@gmail.com}}
            \item {\textit{\url{ http://www.midstorm.org/~fike/weblog}}}
        \end{itemize}
    \end{exampleblock}

\end{frame}



\end{document}
